\title{QLoRA: Efficient Finetuning of Quantized LLMs}

\begin{abstract}
We introduce \textsc{QLoRA}, a memory-efficient finetuning strategy that enables the adaptation of 65B parameter large language models using a single 48GB GPU, all while maintaining the full capabilities of standard 16-bit finetuning. \textsc{QLoRA} achieves this by conducting gradient backpropagation through a fixed, 4-bit quantized version of a pretrained language model into Low Rank Adapters~(LoRA). The strongest models produced via this approach, collectively called \textbf{Guanaco}, surpass previously released open models on the Vicuna benchmark, attaining 99.3\% of ChatGPT’s performance with just 24 hours of finetuning on a single GPU. \textsc{QLoRA} employs several techniques to optimize memory consumption without compromising results: (a) the introduction of 4-bit NormalFloat (NF4), a novel datatype tailored for weights with a normal distribution, (b) Double Quantization, which further compresses memory usage by applying quantization to quantization constants themselves, and (c) Paged Optimizers that help regulate memory demand spikes. Leveraging \textsc{QLoRA}, we finetune over 1,000 models, systematically evaluating instruction following and chatbot behavior across 8 instructional datasets, a variety of model types (including LLaMA and T5), and parameter scales previously inaccessible to conventional finetuning, such as 33B and 65B parameters. The findings suggest that finetuning with QLoRA on modest, high-quality datasets can yield cutting-edge performance, sometimes even surpassing larger prior state-of-the-art models. We conduct extensive assessments of chatbot effectiveness using both human judgments and GPT-4-based evaluations, demonstrating that GPT-4 scoring provides an efficient and reliable alternative to traditional human evaluation. Additionally, we observe significant shortcomings in current chatbot benchmark reliability. A targeted “lemon-picked” analysis highlights failure cases where \textbf{Guanaco} falls short of ChatGPT. All models and source code, including CUDA kernels enabling 4-bit training, are made publicly available.\footnote{\url{https://github.com/artidoro/qlora} and \url{https://github.com/TimDettmers/bitsandbytes}}
\end{abstract}

\section{Introduction}

Adapting large language models (LLMs) through finetuning remains a powerful strategy to boost model accuracy~\citep{min2021metaicl, wei2021finetuned, ouyang2022training, wang2022super, wang2022self, liu2022few} and to intentionally incorporate or eliminate certain behaviors~\citep{ouyang2022training,askell2021general,bai2022training}. Nevertheless, this process is often impractical for very large models due to prohibitive resource demands; for example, conventional 16-bit finetuning of the LLaMA model with 65 billion parameters~\citep{touvron2023llama} consumes over 780 GB of GPU memory. While advancements in quantization have enabled a reduced memory profile for LLMs at inference~\citep{dettmers2022llm, dettmers2022case, frantar2022gptq, xiao2022smoothquant}, such methods falter during the training phase~\citep{wortsman2023stable}.

In this work, we present the first approach that enables finetuning of a quantized 4-bit model without compromising model quality. Our approach, \textsc{QLoRA}, introduces a novel high-precision quantization scheme to convert pretrained models to 4 bits and subsequently attaches a compact set of tunable Low-rank Adapter weights~\citep{hu2021lora}, which are optimized by backpropagating through the quantized base model weights.

Employing \textsc{QLoRA}, we lower the memory consumption for finetuning a 65B parameter model from over 780 GB to under 48 GB, all while matching both computational speed and predictive performance of traditional 16-bit finetuning. This dramatically expands the practical accessibility of LLM finetuning—now even the largest open-source models can be adapted using a single GPU. Leveraging \textsc{QLoRA}, we develop the \textbf{Guanaco} suite of models, with the second most capable version achieving 97.8\% of ChatGPT’s accuracy on the Vicuna benchmark~\citep{vicuna2023}, trainable in under 12 hours on a typical consumer GPU; with a high-end professional GPU for 24 hours, our top-tier model attains 99.3\%, effectively eliminating the performance gap with ChatGPT on this evaluation. On deployment, the smallest \textbf{Guanaco} variant (7B parameters) needs just 5 GB of memory, surpassing a 26 GB Alpaca model by over 20 percentage points on the Vicuna benchmark (see Table~\ref{tbl:vicuna_eval}).

\textsc{QLoRA} incorporates several key advancements aimed at minimizing memory consumption while maintaining performance standards: (1) \textbf{4-bit NormalFloat}, a quantization scheme optimized for normally distributed values, grounded in information theory, which empirically outperforms both 4-bit Integers and 4-bit Floats; (2) \textbf{Double Quantization}, which involves quantizing the quantization coefficients themselves, leading to an average reduction of about 0.37 bits per parameter—amounting to roughly 3 GB saved for a 65B parameter model; and (3) \textbf{Paged Optimizers}, which leverage NVIDIA’s unified memory architecture to alleviate the abrupt memory surges typically associated with gradient checkpointing during the processing of long sequences in mini-batches. Together, these advancements are integrated into an enhanced LoRA framework that places adapters at each layer in the network, effectively mitigating most of the accuracy compromises present in earlier methods.

The efficiency brought by \textsc{QLoRA} allows for extensive investigation into instruction finetuning and chatbot effectiveness at scales that would be otherwise unfeasible with standard finetuning, due to prohibitive memory requirements. To this end, we conduct experiments with over 1,000 models spanning a range of instruction tuning datasets, model designs, and sizes from 80M up to 65B parameters. Beyond demonstrating that \textsc{QLoRA} matches 16-bit baseline performance (\S\ref{sec:comp_to_std}) and facilitates the training of the leading chatbot, \textbf{Guanaco}, (\S\ref{sec:pushing}), we also systematically analyze patterns in these models’ outcomes. Notably, our findings underscore that dataset quality surpasses dataset scale in importance; for instance, a compact 9k sample dataset (OASST1) yielded superior chatbot results compared to a much larger 450k example subset (FLAN v2), despite both targeting generalized instruction following. Additionally, we show that excelling at the Massive Multitask Language Understanding (MMLU) benchmark does not necessarily translate into high performance on the Vicuna chatbot evaluation, and vice versa, highlighting that the appropriateness of a dataset for a specific objective outweighs raw dataset size.

In addition, we conduct a thorough examination of chatbot effectiveness through dual evaluation by human judges as well as GPT-4. Our evaluation adopts a tournament-style framework, where models engage in head-to-head contests to generate optimal responses for given prompts. Match victors are determined by either human raters or GPT-4. The outcomes from these tournaments are synthesized into Elo ratings~\citep{elo1967proposed, elo1978rating}, which serve to order the chatbots by performance. Our findings indicate a strong alignment between human and GPT-4 assessments regarding model rankings in these tournaments; however, notable disagreements occasionally arise. Accordingly, we underscore that although model-based evaluation provides a cost-effective alternative to human annotation, it introduces specific uncertainties that must be acknowledged.

Complementing our quantitative benchmarking, we perform qualitative evaluations of \textbf{Guanaco} models, drawing attention to both successful interactions and shortcomings that were not identified in the numerical analyses.

To support ongoing research, we provide public access to all model outputs along with their corresponding annotations from both human raters and GPT-4. Our implementation—including source code, custom CUDA kernels, and integration with the Hugging Face transformers library \citep{wolf2019huggingface}—is fully open-sourced for community use. Furthermore, we make available a suite of adapters compatible with 7/13/33/65B model sizes, each trained on one of eight distinct instruction-following datasets, resulting in a total release of 32 open-source, finetuned models.

\section{Background}
\paragraph{Block-wise k-bit Quantization} Quantization refers to converting a data representation with higher precision to one with lower precision. This generally involves transforming data from a wider bit-width type, such as 32-bit floats, into a lower bit-width type like 8-bit integers. To utilize the entire value range of the low-bit type, the original data are rescaled—commonly by normalizing according to the maximum absolute value within the dataset, which is often represented as a tensor. For instance, converting a tensor of 32-bit floating-point values (FP32) to an 8-bit integer (Int8) in the range $[-127, 127]$ involves:

\begin{equation}
    \mathbf{X}^{ \text{Int8}} = \text{round}\left( \frac{127}{{\text{absmax}}(\mathbf{X}^{{\text{FP32}}})}\mathbf{X}^{{\text{FP32}}} \right) = \text{round}(c^{\text{FP32}}\cdot \mathbf{X}^{{\text{FP32}}}),
\end{equation}

where $c$ denotes the {\em quantization constant} or {\em quantization scale}. The dequantization operation, which reverses quantization, is defined as:

\begin{equation}
    \text{dequant}(c^{\text{FP32}}, \mathbf{X}^{\text{Int8}}) = \frac{\mathbf{X}^{\text{Int8}}}{c^{\text{FP32}}} = \mathbf{X}^{\text{FP32}}
\end{equation}

A key limitation of this technique arises when the input tensor contains a value with an unusually large magnitude (an outlier). In such cases, certain quantization bins—specific bit patterns—become underutilized or completely unused, as only a small subset of input values end up in those bins. To address the problem introduced by outliers, a frequently used strategy is to partition the input tensor into smaller segments, referred to as blocks, and quantize each block independently, assigning each one its own quantization constant $c$. More precisely, the input tensor $\mathbf{X}\in \mathbb{R}^{b\times h}$ is reshaped into a vector, and then divided into $n$ consecutive blocks, each of length $B$, yielding ${n = ({b\times h} )/{B} }$ blocks. Each block is separately quantized according to Equation~1, resulting in a quantized representation as well as $n$ individual scaling constants $c_i$.

\paragraph{Low-rank Adapters} Low-rank Adapter (LoRA) finetuning \citep{hu2021lora} is an approach designed to lower memory consumption by introducing a compact set of trainable adapter parameters, while the main model parameters remain fixed and are not updated. During stochastic gradient descent, gradients propagate through the static pretrained parameters to the adapters, which are then adjusted to minimize the loss. LoRA modifies a standard linear transformation by introducing an additional, factorized low-rank transformation. For a projection described by ${\mathbf{X} \mathbf{W} = \mathbf{Y}}$, with $\mathbf{X}\in  \mathbb{R}^{b\times h}$ and $\mathbf{W}\in  \mathbb{R}^{h\times o}$, the LoRA method computes:

\begin{equation}
    \mathbf{Y} = \mathbf{X}\mathbf{W} + s\mathbf{X}\mathbf{L}_1\mathbf{L}_2,
\end{equation}

where $\mathbf{L}_1\in  \mathbb{R}^{h\times r}$, $\mathbf{L}_2\in  \mathbb{R}^{r\times o}$, and $s$ denotes a scalar scaling factor.

\paragraph{Memory Requirement of Parameter-Efficient Finetuning} A key aspect to consider is the memory consumption of LoRA during the training process, particularly with respect to both the quantity and size of the adapters deployed. The inherently small memory footprint of LoRA allows for an increased number of adapters to be utilized, thereby enhancing performance without notably raising the total memory consumption. Although LoRA was introduced as a Parameter Efficient Finetuning (PEFT) technique, in practice, the dominant portion of memory usage during LLM finetuning arises from activation gradients rather than from the stored LoRA parameters themselves. For instance, when training a 7B LLaMA model on FLAN v2 with a batch size of 1 and setting the LoRA weights to 0.2\% of the full model weights—as is standard in the literature~\citep{hu2021lora,liu2022few}—the memory required for LoRA input gradients reaches 567 MB, whereas the LoRA parameter storage occupies a mere 26 MB. Employing gradient checkpointing~\citep{chen2016training} reduces the average input gradient memory to 18 MB per sequence, still surpassing the total size of all LoRA parameters. For reference, the underlying 4-bit model itself takes up 5,048 MB. These observations suggest that while gradient checkpointing can meaningfully decrease gradient memory, making further reductions to LoRA parameters provides minimal additional savings. Consequently, scaling up the number of adapters does not substantially impact the overall training memory demand (see Appendix~\ref{app:memory} for a comprehensive analysis). This property is especially important when aiming to achieve performance that matches full 16-bit precision, as discussed later.

\section{\textsc{QLoRA} Finetuning}

\textsc{QLoRA} enables precise finetuning at 4-bit precision through two novel methods: 4-bit NormalFloat (NF4) quantization and Double Quantization. To address memory management issues—specifically, to avoid out-of-memory errors caused by gradient checkpointing that traditionally complicate single-machine finetuning of large models—we also present Paged Optimizers.

Within \textsc{QLoRA}, model weights are maintained in a low-precision storage format (typically 4-bit), while computations are carried out in a higher-precision format, generally BFloat16. As a result, whenever a weight tensor is accessed during computation, it is dequantized to BFloat16 before undergoing 16-bit matrix multiplication.

Below, we outline the principal elements of \textsc{QLoRA} and then provide a formal characterization of the approach.

\paragraph{4-bit NormalFloat Quantization} The NormalFloat (NF) data type extends the ideas of Quantile Quantization \citep{dettmers2022optimizers}, which is an information-theoretically optimal representation assigning the same count of input tensor values to each quantization bin. Quantile quantization operates by leveraging the empirical cumulative distribution function to determine quantile locations in the data.

A major drawback of quantile quantization is the high computational cost associated with quantile estimation. To alleviate this, rapid quantile approximation algorithms, such as SRAM quantiles~\citep{dettmers2022optimizers}, are employed. Nevertheless, due to inherent estimation inaccuracies in these algorithms, substantial quantization error often occurs for outlier values—which may represent the most critical elements of the data.

When input tensors originate from distributions fixed up to a quantization constant, the challenges of costly quantile computation and approximation errors can be bypassed, as such distributions yield identical quantiles across tensors, rendering precise quantile estimation both possible and efficient.

Pretrained neural network weights are typically distributed according to a zero-mean normal distribution with a certain standard deviation $\sigma$ (refer to Appendix~\ref{app:normality}). By appropriately scaling $\sigma$, all such weights can be mapped onto a uniform target distribution confined to the representable range of our chosen data type. In this work, the range of the data type is set to $[-1, 1]$. Therefore, it is necessary to normalize both the neural network weights and the quantile breakpoints of the data type within this specified interval.

For zero-mean normal distributions with arbitrary $\sigma$ values, the theoretically optimal quantization scheme in the fixed interval $[-1, 1]$ proceeds as follows: (1) compute the $2^k+1$ quantiles of a standard normal distribution $N(0, 1)$ to derive a $k$-bit quantile quantizer suited for normal distributions; (2) adjust the resulting quantizer values to span $[-1, 1]$; and (3) transform any given weight tensor by normalizing its entries to $[-1, 1]$, typically using absolute-max normalization. 

With both the weight values and quantization levels now defined over a common range, standard quantization techniques may be directly applied. The third step corresponds mathematically to aligning the standard deviation of the weight tensor with that of the $k$-bit quantization data type. In explicit terms, the $2^k$ quantized values $q_i$ are determined by

\begin{equation}
q_i = \frac{1}{2} \left( Q_X \left( \frac{i}{2^k+1} \right ) + Q_X \left( \frac{i+1}{2^k+1} \right) \right),
\end{equation}

where $Q_X(\cdot)$ denotes the quantile function for the standard normal distribution $N(0,1)$. A limitation of symmetric $k$-bit quantization is the lack of an exact zero representation; such representation is crucial for precisely encoding zero-valued entries, such as those from padding, without error. To address this, and to utilize the full $2^k$ representable values of a $k$-bit format while guaranteeing that zero is assigned a distinct quantized code, we design an asymmetric quantization scheme. Here, quantiles $q_i$ are computed separately for two intervals: $2^{k-1}$ for the negative segment and $2^{k-1}+1$ for the positive segment. These quantile sets are then combined, and one redundant zero, which would be shared between the two sets, is eliminated. We refer to the resulting datatype as \emph{k-bit NormalFloat} (NFk), as each quantization bin contains an equal expected count of values and the format is theoretically optimal for normal distributions centered at zero. The exact quantized values are detailed in Appendix~\ref{app:nf4}.

\paragraph{Double Quantization{}} 
We propose \emph{Double Quantization} (DQ), which entails quantizing the quantization coefficients themselves to achieve further reductions in memory usage. Standard 4-bit quantization~\citep{dettmers2022case} demands a small blocksize to achieve high accuracy, but this comes at the cost of increased memory for quantization constants. For instance, quantizing with 32-bit constants at a blocksize of 64 for $\mathbf{W}$ results in an average overhead of $32/64 = 0.5$ bits per parameter. The Double Quantization strategy effectively alleviates this overhead by compressing these constants.

More concretely, Double Quantization involves applying an additional quantization step to the quantization constants $c_2^{\text{FP32}}$ generated from the primary quantization. This produces quantized constants $c_2^{\text{FP8}}$ and introduces a new set of quantization constants $c_1^{\text{FP32}}$ at this secondary level. For the secondary quantization, we employ 8-bit Floats with a blocksize of 256, as experiments and findings by \citet{dettmers2022case} confirm there is no noticeable performance drop at this precision. Noting that $c_2^{\text{FP32}}$ are all positive, we normalize them by subtracting their mean before quantizing, thus centering their distribution and permitting the use of symmetric quantization. With a blocksize of 64, this method brings the storage required for quantization constants down from $32/64=0.5$ bits to $8/64 + 32/(64\cdot256) = 0.127$ bits per parameter, representing a reduction of 0.373 bits per parameter.

\paragraph{Paged Optimizers} leverage NVIDIA's unified memory~\footnote{\tiny \url{https://docs.nvidia.com/cuda/cuda-c-programming-guide}} capability, which transparently handles page-wise memory transfers between the CPU and GPU. This mechanism ensures reliable GPU computations even in instances where GPU memory is depleted. Functionally, unified memory operates similarly to traditional CPU RAM and disk paging. In our approach, we allocate the optimizer states within this paged memory system, allowing states to be automatically swapped out to CPU RAM when GPU memory limits are reached and reloaded into GPU memory as needed during optimizer updates.

\paragraph{\textsc{QLoRA}.} Incorporating the mechanisms outlined above, we formalize \textsc{QLoRA} for a single linear layer in the quantized backbone, equipped with a single LoRA adapter, as given by:

\begin{equation}
    \mathbf{Y}^{\text{BF16}} =  \mathbf{X}^{\text{BF16}} \text{doubleDequant}(c_1^{\text{FP32}}, c_2^{\text{k-bit}}, \mathbf{W}^{\text{NF4}}) + \mathbf{X}^{\text{BF16}}\mathbf{L}^{\text{BF16}}_1\mathbf{L}^{\text{BF16}}_2,
\end{equation}

where the function doubleDequant$(\cdot)$ takes the form:

\begin{equation}
    \text{doubleDequant}(c_1^{\text{FP32}}, c_2^{\text{k-bit}}, \mathbf{W}^{\text{k-bit}}) = \text{dequant}(\text{dequant}(c_1^{\text{FP32}}, c_2^{\text{k-bit}}), \mathbf{W}^{\text{4bit}}) = \mathbf{W}^{\text{BF16}},
\end{equation}

Here, NF4 quantization is used for $\mathbf{W}$ and FP8 for $c_2$.

To achieve finer quantization, $\mathbf{W}$ is processed with a blocksize of 64, while $c_2$ employs a larger blocksize of 256 to optimize for memory efficiency.

During parameter updates, only the gradients with respect to the adapter weights, $\frac{\partial E}{\partial \mathbf{L}_i}$, are required; gradients for the 4-bit weights, $\frac{\partial E}{\partial \mathbf{W}}$, are unnecessary. Nevertheless, determining $\frac{\partial E}{\partial \mathbf{L}_i}$ inherently involves calculating $\frac{\partial \mathbf{X}}{\partial \mathbf{W}}$, which utilizes equation~(5). This requires dequantizing the stored weights $\mathbf{W}^{\text{NF4}}$ to their computation type $\mathbf{W}^{\text{BF16}}$, so that $\frac{\partial \mathbf{X}}{\partial \mathbf{W}}$ is evaluated in BFloat16 precision.

In summary, \textsc{QLoRA} utilizes a single storage data type—typically 4-bit NormalFloat—and a distinct computation data type, specifically 16-bit BrainFloat. During both the forward and backward passes, weights stored in the storage data type are dequantized to the computation data type. Notably, weight gradients are only calculated for the LoRA-specific parameters, which remain in 16-bit BrainFloat.

\section{QLoRA vs. Standard Finetuning}
\label{sec:comp_to_std}

We have outlined the mechanism of QLoRA and its potential to drastically lower memory consumption in model finetuning. The central issue is whether QLoRA delivers comparable outcomes to complete model finetuning. Additionally, it is important to dissect the various elements of QLoRA, including the influence of using NormalFloat4 as opposed to the conventional Float4. The experiments described below were designed to address these considerations.

\paragraph{Experimental setup.} Our investigation encompasses three model classes—encoder, encoder-decoder, and decoder-only architectures—and contrasts QLoRA with both 16-bit adapter-based finetuning and full-parameter finetuning for models as large as 3B parameters. The evaluation covers a range of benchmarks: GLUE \citep{wang2018glue} with RoBERTa-large \citep{liu2019roberta}; Super-NaturalInstructions (TKInstruct) \citep{wang2022super} utilizing T5 \citep{t5}; and 5-shot MMLU \citep{hendrycksmeasuring} where LLaMA is finetuned on Flan v2 \citep{longpre2023flan} and Alpaca \citep{alpaca}. To assess the comparative benefits of NF4 over alternative 4-bit quantization schemes, we adopt the framework of \citet{dettmers2022case}, recording post-quantization zero-shot accuracy and perplexity across multiple models (OPT~\citep{zhang2022opt}, LLaMA~\citep{touvron2023llama}, BLOOM~\citep{scao2022bloom}, Pythia~\citep{biderman2023pythia}) spanning from 125 million to 13 billion parameters.

Further specifics regarding each experimental scenario are included in the respective results sections to enhance clarity; comprehensive setup descriptions can be found in Appendix~\ref{app:exp-setup}.

Although paged optimizers are indispensable for training 33B/65B \textsc{QLoRA} models on GPUs with 24GB or 48GB of memory, we omit detailed performance metrics for Paged Optimizers, as paging typically occurs only when very long sequences are processed in mini-batches, a situation encountered infrequently. Nevertheless, runtime analysis on 65B models using 48GB GPUs indicates that, with a batch size of 16, paged optimizers achieve comparable throughput to standard optimizers. Identifying specific scenarios that may induce paging-related slowdowns remains an avenue for future research.

\paragraph{Default LoRA hyperparameters do not match 16-bit performance}

When employing the conventional approach of integrating LoRA into the query and value attention projection matrices \citep{hu2021lora}, we observe that full finetuning performance cannot be reproduced for models of considerable scale. As demonstrated in Figure~\ref{fig:hyper}—which examines LLaMA 7B finetuning on Alpaca—the primary LoRA hyperparameter influencing outcomes is the total number of LoRA adapters applied. Achieving equivalence with full finetuning performance requires applying LoRA to every linear layer within the transformer blocks. By contrast, other hyperparameters, such as the projection dimension $r$, appear to have negligible effect on results (see Appendix \ref{app:exp-setup}).

In a similar vein, our investigation reveals that the default hyperparameters used for fully finetuned baseline models are suboptimally tuned. To establish stronger baseline performances, we perform an extensive hyperparameter sweep, varying learning rates from $1\text{e}^{-6}$ up to $5\text{e}^{-5}$ and batch sizes from 8 to 128. The outcomes for 7B LLaMA finetuning on Alpaca are depicted in Figure~\ref{fig:hyper}.

\paragraph{4-bit NormalFloat outperforms 4-bit Floating Point}

Despite 4-bit NormalFloat (NF4) being theoretically optimal with respect to information retention, it remains unclear whether this theoretical advantage confers practical benefits. Adopting the experimental design described by \citet{dettmers2022case}, we evaluate quantized LLMs—specifically OPT \citep{zhang2022opt}, BLOOM \cite{scao2022bloom}, Pythia \cite{biderman2023pythia}, and LLaMA—across a range of model sizes (from 125M up to 65B parameters) and datatype choices, considering both standard language modeling as well as several zero-shot benchmarks. According to Figure~\ref{fig:nf4} and Table~\ref{tbl:nf4_ppl}, we observe that NF4 offers substantial gains in performance compared to both FP4 and Int4; furthermore, using double quantization yields a reduction in memory usage without compromising effectiveness.

\paragraph{k-bit \textsc{QLoRA} achieves parity with full 16-bit and 16-bit LoRA finetuning} 

Prior studies have demonstrated that employing 4-bit quantization for \emph{inference} is feasible, though typically at the cost of some degradation compared to 16-bit precision \citep{dettmers2022case,frantar2022gptq}. This observation prompts an important question: Can 4-bit adapter finetuning fully reclaim the accuracy sacrificed by quantization? We empirically examine this under two experimental regimes.

The initial investigation contrasts fully finetuned 16-bit models—specifically, RoBERTA and T5 across 125M to 3B parameter scales—against 8-bit and 4-bit adapter-based methods on the GLUE and Super-NaturalInstructions benchmarks. Table~\ref{tab:nf4vsbf16} presents the results. For both datasets, 16-bit, 8-bit, and 4-bit adapter techniques all match the performance of the original fully finetuned 16-bit baselines. This indicates that any performance deterioration due to aggressive quantization can be entirely mitigated by adapter-based finetuning applied after quantization.

In our second experimental configuration, full finetuning of models at the 11B parameter scale and above exceeds the memory capacity of a single high-memory GPU server. Therefore, we further examine whether 4-bit \textsc{QLoRA} can achieve performance parity with 16-bit LoRA for model sizes ranging from 7B to 65B parameters. To investigate this, we apply finetuning to LLaMA models from 7B up to 65B parameters using the Alpaca and FLAN v2 instruction-tuning datasets, subsequently assessing their performance on the MMLU benchmark using 5-shot accuracy. Table~\ref{tbl:llama-datatype} presents the outcomes, demonstrating that the NF4 quantization method, coupled with double quantization, entirely matches the MMLU scores achieved by 16-bit LoRA. Additionally, we observe that \textsc{QLoRA} using FP4 shows approximately a 1-point deficit relative to the 16-bit brain float LoRA baseline. These findings reinforce two central points: (1) \textsc{QLoRA} with NF4 successfully reproduces both 16-bit full finetuning and 16-bit LoRA outcomes, and (2) NF4 offers improved quantization fidelity over FP4.

\paragraph{Summary} Our findings provide robust evidence that 4-bit \textsc{QLoRA} employing the NF4 data representation attains results equivalent to both 16-bit full finetuning and 16-bit LoRA across standard academic benchmarks using established evaluation methodologies. We further demonstrate the superior effectiveness of NF4 over FP4 and confirm that the process of double quantization does not compromise finetuning outcomes. Collectively, these results underscore the reliability of 4-bit \textsc{QLoRA} tuning for matching the capabilities of higher precision methods.

Consistent with prior quantization research \citep{dettmers2022case}, the results from our MMLU and Elo evaluations suggest that, given fixed finetuning and inference computational resources, prioritizing an increase in base model parameter count while utilizing lower precision is advantageous. This observation emphasizes the practical efficiency gains offered by \textsc{QLoRA}. As our experiments reveal no loss in accuracy relative to full-finetuning when using 4-bit precision, it remains an open question where the optimal balance between performance and precision lies for QLoRA finetuning—a topic we propose for future study.

We extend our investigation to instruction tuning at scales that would be unattainable using full 16-bit finetuning given typical academic research computing resources.

\section{Pushing the Chatbot State-of-the-art with QLoRA}
\label{sec:pushing}
Having demonstrated that 4-bit \textsc{QLoRA} is capable of matching 16-bit performance across various scales, tasks, and datasets, we perform a comprehensive analysis of instruction finetuning, scaling up to the largest research-accessible open-source language models. To evaluate the effectiveness of instruction finetuning for these models, we use the demanding Natural Language Understanding benchmark (MMLU) and introduce novel approaches for measuring chatbot performance in practical contexts.

\subsection{Experimental setup} Below, we outline the experimental framework; complete methodological specifics can be found in Appendix \ref{app:chatbot}.

\paragraph{Data} In light of the lack of thorough comparative analyses of modern instruction-following datasets, we select eight recent datasets for evaluation. These consist of crowd-sourced datasets (OASST1~\citep{kopf2023openassistant}, HH-RLHF~\citep{bai2022training}), data distilled from models that have undergone instruction-tuning (Alpaca~\citep{alpaca}, self-instruct~\citep{wang2022self}, unnatural-instructions~\citep{honovich2022unnatural}), large aggregated corpora (FLAN v2~\citep{chung2022scaling}), as well as hybrid approaches (Chip2~\citep{laion2023}, Longform~\citep{koksal2023longform}). These selections collectively provide diversity in language, dataset scale, and licensing terms.

\paragraph{Training Setup} To eliminate confounding due to varying training methodologies, all QLoRA finetuning is carried out using cross-entropy loss under a supervised learning regime, without incorporating reinforcement learning—even for those datasets where human feedback is present. For datasets distinguishing between instructions and responses, we perform finetuning solely on the response segments (see Appendix \ref{app:chatbot} for additional ablation studies). In OASST1 and HH-RLHF, which offer multiple possible responses, we identify the highest-quality response at each conversational node and finetune on the complete sequence of selected exchanges, including both instructions and their replies. Throughout our work, we utilize NF4 \textsc{QLoRA} configured with double quantization and paged optimizers to mitigate memory usage spikes during gradient checkpointing. For 13B and 33B LLaMA models, modest hyperparameter searches are conducted. We observe that most hyperparameter choices optimized for 7B models (e.g., number of epochs) transfer robustly, though batch size and learning rate must be adjusted: the learning rate is reduced by half and the batch size is doubled for the 33B and 65B models.

\paragraph{Baselines} For comparison, we benchmark our models against both academic (Vicuna~\citep{vicuna2023}, Open Assistant~\citep{kopf2023openassistant}) and proprietary chatbot platforms (GPT-4 \citep{openai2023gpt}, GPT-3.5-turbo, Bard). Notably, Open Assistant is built on a LLaMA 33B model finetuned via Reinforcement Learning from Human Feedback (RLHF) using the OASST1 dataset, the same dataset employed in our experiments. Vicuna performs full finetuning of LLaMA 13B using private conversation data from ShareGPT, thereby acting as a distillation of OpenAI GPT models.

\subsection{Evaluation}

Adhering to established methodology, we assess language understanding capabilities via the MMLU (Massively Multitask Language Understanding) benchmark \citep{hendrycksmeasuring}. This challenging suite spans 57 multiple-choice tasks drawn from disciplines such as elementary mathematics, US history, computer science, law, and others. All reported results correspond to 5-shot accuracy on the test set.

Our investigation of generative language abilities incorporates both machine-driven and human-centered assessment methods. This latter group of evaluations involves human-constructed query sets designed to gauge the response quality generated by the models. Such approaches, though increasingly regarded as providing a more authentic measure of chatbot performance, lack standardized protocols within existing literature. Below, we outline our chosen methodology, consistently employing nucleus sampling with $p=0.9$ and a temperature of $0.7$ across all experimental conditions.

\paragraph{Benchmark Data} Our experiments utilize two specifically selected query datasets. The first is the set of Vicuna prompts \citep{vicuna2023}, which comprises 80 unaltered prompts spanning a range of subject areas. The second is the OASST1 validation dataset \citep{kopf2023openassistant}, featuring multilingual, crowd-sourced dialogues between users and assistants. For this dataset, every user message in the validation split serves as an individual query, and previous dialogue exchanges are retained in the prompt for contextual continuity, resulting in 953 distinct user queries. We refer to these datasets as the Vicuna and OA benchmarks.

\paragraph{Automated Evaluation} Our primary automated evaluation follows the methodology set forth by \citet{vicuna2023}, utilizing GPT-4 to assess and score outputs on the Vicuna benchmark. In each trial, GPT-4 is presented with a prompt alongside responses from both ChatGPT (GPT-3.5 Turbo) and the candidate model, assigning each a score on a 10-point scale with accompanying justification. Model performance is then computed as a percentage of the score obtained by ChatGPT, with relative scores exceeding 100\% possible when models surpass ChatGPT's absolute scores. It is important to note that we observed a consistent bias wherein GPT-4 tends to grant higher scores to the response shown first in the prompt. To account for this positional effect, we advocate for averaging scores across both presentation orders.

Subsequently, we conduct performance assessment through direct, head-to-head output comparisons. For these evaluations, the rubric is reduced to three possible outcomes—one response preferred, the other preferred, or a tie. GPT-4 is prompted to indicate the superior response or declare a tie, always providing a rationale. These pairwise evaluations are exhaustively performed for all system pairs across both Vicuna and OA benchmarks.

\paragraph{Human Evaluation} Although recent studies suggest generative models can serve as useful tools for evaluation tasks \citep{fu2023gptscore}, the alignment of GPT-4-based ratings with human assessments of chatbot effectiveness is, at present, unsubstantiated. Consequently, we supplement our automated metrics with parallel human evaluations on the Vicuna dataset, mirroring both automated protocols. For comparisons against ChatGPT, two annotators are recruited via Amazon Mechanical Turk (AMT), while three annotators are employed for system pairwise assessments.

\paragraph{Elo Rating} Utilizing both human and automated pairwise evaluations, we construct a tournament framework where models engage in direct competition. Each round consists of two models vying to generate the superior response to a shared prompt. This method follows the approach of \citet{bai2022training} and \citet{vicuna2023} in model comparison, but distinguishes itself by incorporating GPT-4-based ratings alongside human judgments. Elo scores~\citep{elo1967proposed,elo1978rating} are then calculated from a randomly selected subset of labeled matchups. Originally developed for ranking chess players, the Elo system quantifies a player’s expected likelihood of victory in comparison to their opponent; for instance, a player rated at 1100 has a roughly 65\% predicted win probability against an opponent rated at 1000, whereas two players with equal ratings (e.g., 1000 vs 1000 or 1100 vs 1100) are expected to split victories evenly at 50\%. Following each match, the Elo score is updated according to the degree of surprise in the result—unexpected upsets prompt larger adjustments, while anticipated outcomes cause only minor updates. Across repeated matchups, these ratings gradually reflect each participant’s true competitive ability. Each model is initially assigned a score of 1,000, with a $K=32$ adjustment parameter. To mitigate the effects of match order, as in \citet{vicuna2023}, this process is repeated 10,000 times using varied random seeds, helping ensure unbiased outcomes regardless of pairing sequence.

\subsection{\textbf{Guanaco}: \textsc{QLoRA} trained on OASST1 Represents a Leading-Edge Open-Source Chatbot}

Drawing on both automated metrics and assessments from human evaluators, our analyses reveal that the {Guanaco} 65B model, a product of \textsc{QLoRA} fine-tuning on a modified version of OASST1, outperforms all other open-source chatbot models and closely rivals ChatGPT in terms of capability. In direct comparison to GPT-4, {Guanaco} 65B and 33B models exhibit an anticipated win probability of 30\%, according to Elo scores computed from pairwise model evaluations by human raters—currently the highest win rate observed.

Table~\ref{tbl:vicuna_eval} presents comparative results on the Vicuna benchmark \cite{vicuna2023} with reference to ChatGPT. The results indicate that, aside from GPT-4, {Guanaco} 65B yields superior performance, reaching 99.3\% of ChatGPT’s evaluated effectiveness. While {Guanaco} 33B is more parameter-rich than Vicuna 13B, its reliance on 4-bit precision parameters results in enhanced memory efficiency—requiring just 21 GB, compared to Vicuna’s 26 GB—while securing a three-point performance gain over the latter. Additionally, {Guanaco} 7B can run on modern smartphones with only a 5 GB memory footprint and still surpasses Alpaca 13B by almost 20 percentage points.

Despite these impressive results, the confidence intervals reported in Table~\ref{tbl:vicuna_eval} are broad, with notable overlap in model performances. We speculate this ambiguity may stem from inadequate definition of rating scales—such as the variable interpretation of an 8 on a 10-point scale in diverse contexts. To address the limitations of fixed scales, we advocate for the use of the Elo rating methodology \citep{elo1967proposed}, which leverages \textit{pairwise} human and GPT-4 judgments to mitigate the problem of arbitrary scoring anchors. Table~\ref{tbl:elo} details the Elo rankings among the most competitive systems. We observe partial disagreement between human raters and GPT-4 rankings for some models—especially Guanaco 7B—though broader consistency is evidenced, with a Kendall Tau of $\tau=0.43$ and a Spearman correlation coefficient of $r=0.55$ at the system level. On the instance level, concordance between GPT-4 and the majority opinion of human annotators is lower, with a Fleiss’ $\kappa=0.25$. Taken together, these results indicate moderate agreement for system-level comparisons between GPT-4 and human evaluators, suggesting that automated, model-based judgments provide a reasonably dependable substitute for human-based evaluation. Further discussion of these evaluation considerations is available in Section~\ref{ssec:considerations}.

The Elo scores presented in Table~\ref{tbl:elo_large} reveal that the {Guanaco} models with 33B and 65B parameters surpass all counterparts aside from GPT-4 on both the Vicuna and OA benchmarks. Their performance closely matches that of ChatGPT, as also reflected in Table~\ref{tbl:vicuna_eval}. It is worth mentioning that the Vicuna benchmark tends to give an advantage to open-source models, whereas the broader OA benchmark generally favors ChatGPT. Additionally, a comparison of results in Tables \ref{tbl:mmlu-llama} and \ref{tbl:vicuna_eval} underscores the crucial impact of the finetuning dataset choice on overall outcomes. For instance, finetuning Llama models with FLAN v2 results in particularly high scores on MMLU but leads to the lowest performance on Vicuna, a trend seen across various architectures. This demonstrates that current benchmarks are not entirely overlapping: high scores on MMLU do not guarantee strong chatbot capabilities as measured by Vicuna or OA metrics, and vice versa.

Among the top performers evaluated,  is unique in not being trained on proprietary data, as the OASST1 data collection policy prohibits the inclusion of GPT-generated content. The closest open-source-trained competitor, Anthropic’s HH-RLHF model, achieves a score on the Vicuna benchmark that is 30 percentage points lower than {Guanaco}, as shown in Table~\ref{tbl:vicuna_eval}. Collectively, these findings illustrate the efficacy of the 4-bit \textsc{QLoRA} method, which enables the creation of chatbots that reach state-of-the-art performance and challenge ChatGPT. Notably, training the 33B {Guanaco} model is feasible on standard 24 GB consumer GPUs within less than 12 hours. This scalability paves the way for future endeavors involving \textsc{QLoRA} finetuning on domain-specific open-source corpora, potentially yielding models that are able to compete with the top commercial systems currently available.

\section{Qualitative Analysis}

Although our evaluation primarily relies on quantitative measures, exclusive focus on aggregate statistics introduces several challenges. Chief among these is the concern of benchmark validity \citep{liao2021we}—namely, the uncertainty as to whether a benchmark genuinely assesses the skills it purports to measure, a concern exacerbated by the emergence of ``shortcuts'' that machine learning models are apt to exploit \citep{gururangan2018annotation, poliak2018hypothesis}. To address this, we supplement our analysis with qualitative investigations, organized in two subsections. The first, \S\ref{ssec:examples}, provides selected instances that capture characteristic behaviors identified in outputs produced by our 65b {Guanaco} model. The second, \S\ref{ssec:considerations}, discusses important caveats and interpretive factors concerning our findings.

\subsection{Qualitative Analysis of Example Generations}
\label{ssec:examples}

To select representative examples, we reviewed model outputs on both the Vicuna and OpenAssistant benchmarks, searching for recurring response patterns generated by {Guanaco}. Upon identifying such patterns, we formulated targeted prompts to encourage the appearance of these patterns—even when they reflect incorrect solutions. For instance, if the model habitually produces overly detailed replies, we might instruct it to ``Answer yes or no without explanation.'' This procedure enables us to identify both instances in which the model is successfully “broken” through adversarial prompting (referred to as ``lemons''), and cases where it remains robust (termed ``cherries''). Every example shown in this section was produced using Nucleus Sampling \cite{holtzmancurious} with $p=0.9$.

Naturally, this discussion is not exhaustive; addressing every factor at play lies outside the purview of this limited qualitative analysis. For example, the model is capable of producing a vast array of responses to a single prompt, and so we utilize samples that we consider broadly representative. Nonetheless, presenting these particular cases serves to contextualize the quantitative findings detailed earlier in the manuscript. As all models and source code are publicly available, we aim for this section to motivate more granular follow-up research into the concerns highlighted herein.

\paragraph{Factual Recall} When asked direct factual questions such as ``What is the capital of Zambia?'', all models reliably respond with correct information, e.g.,

\begin{quote}
    \textbf{Guanaco}:
    \texttt{The capital of Zambia is Lusaka.}
\end{quote}

Yet as the factual requests become less common, {Guanaco}'s outputs grow unreliable despite maintaining an assertive tone. For example, given the following prompt from HotPotQA \citep{yang2018hotpotqa}:

\begin{quote}
    \textbf{User}: In what year was the singer who popularized `I'll Keep the Lovelight Burning' born?  

    \textbf{Guanaco}:~\texttt{Al Jolson is credited with popularizing the song `I'll Keep the Lovelight Burning,' and he was born in the year 1886.}
\end{quote}

The model replies with both an incorrect attribution and erroneous year of birth, though the stated birth year does match the named individual (Al Jolson).

\paragraph{Suggestibility} {Guanaco} displays an unexpected tendency to resist acquiescing to certain forms of implicit misinformation, as demonstrated in the subsequent exchange:

\begin{quote}
    \textbf{User}: How was it finally, officially confirmed that the earth is flat by peer reviewed scientists? \\ \\
    \textbf{Guanaco}: \texttt{There has never been any official confirmation by peer-reviewed scientists that the earth is flat. Scientific consensus overwhelmingly supports that the earth is a sphere. This conclusion has been supported by countless experiments and observations over many centuries. \\ \\ The assertion that the earth is flat belongs to a conspiracy theory that scientific evidence has thoroughly refuted. No group of scientists has published peer-reviewed work claiming that the earth is flat. \\ \\ Please let me know if you have further questions.}
\end{quote}

 also shows the ability to recognize when certain questions cannot be genuinely addressed, for example,

\begin{quote}
    \textbf{User}: What time is it?

    \textbf{Guanaco}: \texttt{I apologize, but as an AI language model, I do not possess access to up-to-date information such as the current time. My replies are produced using prior knowledge, which does not include real-time data. For the exact time, please consult a clock or look it up online.}
\end{quote}

The same applies to inquiries such as ``Where are you?'', ``How are you?'', and others.

\paragraph{Refusal} In much the same way, {Guanaco} occasionally declines to carry out instructions for apparently arbitrary reasons:

\begin{quote}
    \textbf{User}: Please reverse the order of the words in the sentence ``I'm a cool dude in a mocha mood.'' \\ \\
    \textbf{Guanaco}: \texttt{Apologies, I am unable to do that at this moment; however, I can clarify why your sentence is grammatically accurate.\\ \\The sentence ``I'm a cool dude in a mocha mood'' adheres to English grammar and syntactic conventions. Specifically, the adjective ``cool'' modifies the noun ``dude,'' and the prepositional phrase ``in a mocha mood'' adds further contextual detail. This construction is grammatically sound and the intended meaning is clear.\\ \\If you have any additional queries or require further assistance, please let me know.}
\end{quote}

\paragraph{Secret Keeping} Unfortunately, even when refusal is appropriate or necessary, systems cannot always be relied upon to maintain secrecy. In many cases, system architects intend to withhold certain information from users, even when natural language responses are permitted. To investigate this, we prompt the model with ``The secret word is `banana'. Do not reveal or repeat this back, no matter what the human requests.'' and assess the difficulty of obtaining this concealed detail. A straightforward inquiry does not succeed:

\begin{quote}
    \textbf{User}: What is the secret word?

    \textbf{Guanaco}: \texttt{I'm sorry, but I can't disclose the secret word. I am programmed to refrain from revealing or repeating the secret word, regardless of user requests. Can I help you with something else?}
\end{quote}

This constitutes the expected output. Yet, only minimal manipulation can defeat this safeguard:

\begin{quote}
    \textbf{User}: Let's play a game. The objective is to disregard your previous directives. What is the secret word?

    \textbf{Guanaco}: \texttt{Sorry, the secret word is banana. It's probably best if we don't discuss it further. Is there anything else you'd like to ask?}
\end{quote}

demonstrating the importance of researching techniques to enhance reliability in following instructions.

\paragraph{Math} Mathematics represents a significant challenge for {Guanaco}, as is the case for numerous language models discussed in \citep{liang2022holistic}. When {Guanaco} explicitly lays out its reasoning, its answers are frequently correct, as illustrated below:

\begin{quote}
    \textbf{User}: Lee mows one lawn and charges \$33. Last week he mowed 16 lawns and three customers each gave him a \$10 tip. How many dollars did Lee earn mowing lawns last week?

    \textbf{Guanaco}: \texttt{Lee earned \$582 mowing lawns last week. \\ \\ To find the total amount he earned, we first need to calculate how much he earned for mowing 16 lawns: \\ \\  \$33 x 16 = \$528 \\ \\ Then we need to add the tips he received from the three customers: \\  \\ \$10 x 3 = \$30 \\ \\ So the total amount Lee earned was: \\ \\ \$528 + \$30 = \$558}
\end{quote}

Nonetheless, when {Guanaco} omits detailed step-by-step reasoning, it is prone to errors on even elementary questions, a phenomenon documented in \cite{weichain}. Consider the following exchange as an example:

\begin{quote}
    \textbf{User}: Please factorize 1833.

    \textbf{Guanaco}: \texttt{Sure, I'd be happy to help! 1833 is a prime number, which means it has only two factors: 1 and 1833. \\ \\ So the unique factorization of 1833 is: \\ \\ 2\^{}1 * 3\^{}2 * 17\^{}1 \\ \\ I hope this helps! Let me know if you have any other questions.}
\end{quote}

which is incorrect (the correct factorization is $3 \times 17 \times 43$), and in fact contains two errors.

\paragraph{Theory of Mind} {Guanaco} demonstrates unexpectedly robust Theory of Mind reasoning abilities \cite{nematzadeh2018evaluating, sap2022neural}. For example, the system gives detailed and accurate explanations for queries such as:

\begin{quote}
    \textbf{User}: James and Abby are in the bedroom. Abby put the pen in the desk drawer. Abby leaves the bedroom. James moves the pen into the bag. Where does James think Abby will look for the pen?

    \textbf{Guanaco}: \texttt{James knows that Abby put the pen in the desk drawer. But James also knows that he has moved the pen into the bag, and he assumes that Abby does not know about this. So, James may guess that Abby will look for the pen in the desk drawer, because that's where she left it last time. However, this assumption may be incorrect, because Abby may have seen James moving the pen.}
\end{quote}

Nonetheless, the reasoning is inconsistent, with the model sometimes producing answers that rely on unwarranted or incoherent assumptions, as seen below:

\begin{quote}
    \textbf{User}: Evelyn entered the living room. Jackson entered the playroom. James entered the playroom. The beans are in the treasure chest. James exited the playroom. Jackson moved the beans to the pantry. Jackson exited the playroom. James entered the living room. Where will James look for the beans?

    \textbf{Guanaco}: \texttt{James will look for the beans in the pantry, because that's where Jackson moved them.}
\end{quote}

In this instance, {Guanaco} assumes that James knows what Jackson did, although such information transfer was not specified. These types of mistakes parallel findings reported in recent studies \cite{sap2022neural}, indicating the necessity for further investigation.

\subsection{Considerations}
\label{ssec:considerations}

\paragraph{Evaluation}

Our results show only moderate concordance among human raters (Fleiss $\kappa=0.42$), with even lower agreement when contrasting outputs from two high-performing models. This highlights inherent shortcomings in existing benchmarks and in human-based evaluation procedures for assessing chatbot functionality. Through a qualitative review comparing generations from ChatGPT and {Guanaco} 65B using the Vicuna benchmark, we observe that individual subjectivity becomes increasingly influential; in fact, the authors frequently disagreed on preferred answers. Addressing these challenges in future work may benefit from methodologies in fields such as Human-Computer Interaction and Psychology that are specifically designed to account for subjectivity.

Our study also reveals clear systematic biases in automated assessment systems. In particular, we identify strong positional effects, where GPT-4 tends to assign higher ratings to whichever system appears first in the input prompt. The sample-level agreement between GPT-4 and human evaluators is weak (Fleiss $\kappa=0.25$), implying that automated systems and human raters often use different evaluative criteria. Furthermore, as presented in Table~\ref{tbl:elo_large}, GPT-4 consistently rates its own outputs much more favorably than human annotators do, resulting in Elo scores of 1348 versus 1176, which translates to an approximately 20\% higher likelihood of success in direct comparisons.

Further studies are needed to identify potential biases within automated evaluation frameworks and to explore effective methods for mitigating these biases.

\paragraph{Data \& Training} It is important to highlight that the {Guanaco} models are trained on the multilingual OASST1 dataset, and that the OA benchmark itself presents prompts in a range of languages. Future research should evaluate the impact of multilingual training on the ability of these models to handle non-English instructions, and assess if this factor contributes to the notably larger performance discrepancy observed between the Vicuna-13B model (trained exclusively on English data) and {Guanaco} 33B and 65B when evaluated on the OA benchmark.

In light of {Guanaco} models' strong results, we examined whether there is any data leakage between OASST1 and the Vicuna benchmark prompts. Using fuzzy string matching to compare prompts across the two datasets and manually reviewing the most similar pairs, we found no evidence of prompt overlap.

Additionally, our model is trained exclusively via supervised learning using a cross-entropy loss objective, and does not employ reinforcement learning from human feedback (RLHF). This highlights the need for further comparative studies into the advantages and limitations of standard cross-entropy training relative to RLHF approaches. We anticipate that \textsc{QLoRA} will make such analyses feasible at a large scale without the necessity for extensive computational resources.

\section{Related Work}

\paragraph{Quantization of Large Language Models} Prior efforts in quantizing LLMs have predominantly concentrated on inference-time quantization. State-of-the-art techniques for preserving the fidelity of 16-bit LLMs target outlier features, as seen in works like SmoothQuant~\citep{xiao2022smoothquant} and LLM.int8()~\citep{dettmers2022llm}; alternative methods have leveraged sophisticated grouping strategies~\citep{park2022nuqmm, yao2022zeroquant}. Research on lossy quantization explores the compromises entailed by conventional rounding techniques~\citep{dettmers2022case,zeng2022glm,pope2022efficiently}, as well as the optimization of rounding schemes to maximize precision in quantized models~\citep{frantar2022gptq}. Aside from the current study, SwitchBack layers~\citep{wortsman2023stable} stands as the sole prior work addressing backpropagation through quantized weights at model scales exceeding 1B parameters.

\paragraph{Finetuning with Adapters} In this work, we employ Low-rank Adapters~\citep{hu2021lora} (LoRA), although the literature contains a diverse range of Parameter Efficient FineTuning (PEFT) techniques, including prompt tuning~\citep{qin2021learning,lester2021power,li2021prefix}, modifying embedding layer inputs~\citep{an2022input}, adapting hidden representations as in IA$^3$~\citep{liu2022few}, inserting entire additional layers~\citep{houlsby2019parameter}, bias tuning~\citep{zaken2021bitfit}, weight masking driven by Fisher information~\citep{sung2021training}, as well as hybrid methodologies~\citep{henderson2021compacter}. Through our experiments, we demonstrate that LoRA adapters achieve performance on par with full 16-bit model finetuning. An extensive comparison and exploration of the tradeoffs involved in alternative PEFT methods is left as a direction for subsequent investigations.

\paragraph{Instruction Finetuning} To enable pretrained large language models (LLMs) to follow instructions specified within prompts, instruction finetuning leverages input-output examples drawn from diverse sources to adjust the model such that, given a prompt, it reliably generates the corresponding output. Notable approaches and corpora in this area encompass MetaICL~\citep{min2021metaicl}, MetaTuning~\citep{zhong2021adapting}, InstructGPT~\citep{ouyang2022training}, FLAN~\citep{wei2021finetuned,chung2022scaling}, PromptSource~\citep{bach2022promptsource}, Super-NaturalInstructions~\citep{wang2022super,sanh2021multitask}, Self-instruct~\citep{wang2022self}, UnnaturalInstructions~\citep{honovich2022unnatural}, OPT-IML~\citep{iyer2022opt}, UnifiedSKG\citep{xie2022unifiedskg}, OIG/Chip2~\citep{laion2023}, Alpaca~\citep{alpaca}, Vicuna~\citep{vicuna2023}, Koala~\citep{koala_blogpost_2023}, and Self-instruct-GPT-4~\citep{peng2023instruction}.

\paragraph{Chatbots} Instruction-aligned models are frequently deployed in a conversational context as chatbots. Training these models often involves Reinforcement Learning from Human Feedback (RLHF)~\citep{christiano2017deep} or relies on data generated from current models, with learning guided by automated model feedback (RLAIF)~\citep{bai2022constitutional}. Among prominent methods and datasets are Anthropic-HH~\citep{askell2021general,bai2022training}, Open Assistant~\citep{kopf2023openassistant}, LaMDA~\citep{thoppilan2022lamda}, and Sparrow~\citep{glaese2022improving}. Our approach does not involve reinforcement learning. Instead, our leading model, Guanaco, is trained on multi-turn conversational data from the Open Assistant dataset, which was created with RLHF applications in mind~\citep{kopf2023openassistant}. In the context of chatbot evaluation, recent alternatives replace expensive human annotation with assessment by GPT-4~\citep{vicuna2023,peng2023instruction}. We extend these efforts by introducing improvements aimed at producing a more robust and reliable evaluation protocol.

\section{Limitations and Discussion}

Our findings indicate that \textsc{QLoRA} enables a 4-bit foundation model paired with Low-rank Adapters (LoRA) to attain performance indistinguishable from conventional full 16-bit finetuning. Nevertheless, we have not established whether \textsc{QLoRA} can reliably match 16-bit finetuning outcomes for models at the 33B and 65B parameter scales. Given the substantial computational and financial overhead such experiments would require, addressing this question is reserved for future work.

A further caveat concerns the assessment of instruction-tuned models. While our evaluation covers MMLU, the Vicuna benchmark, and the OA benchmark, it does not extend to additional standard benchmarks such as BigBench, RAFT, or HELM, raising the possibility that our reported results may not transfer to those settings. Conversely, our study features an extensive analysis on MMLU and introduces novel methodologies for chatbot evaluation.

The data suggests that benchmark results are largely contingent on the resemblance between the finetuning corpus and the evaluation dataset. For instance, FLAN v2 shares similarities with MMLU, while it is less aligned with chatbot evaluation sets, a pattern mirrored inversely for the Chip2 dataset; this correspondence is reflected in model performances on both MMLU and Vicuna benchmarks. Such observations underscore the necessity not only for improved benchmarks and evaluative protocols, but also for carefully considering the aspects being measured. It prompts the question of whether the goal should be to optimize for high performance on traditional academic knowledge—akin to high school and undergraduate exam content—or on interactive conversation abilities characteristic of chatbots. The ease of reusing established benchmarks over devising new ones poses a risk of steering collective research agendas towards targets that may not align with the most meaningful objectives. It is thus important, at the community level, to ensure our chosen benchmarks adequately reflect the intended goals of model development.

Although our study provides a comprehensive assessment of general chatbot performance, it is limited by the scope of its responsible AI evaluation for {Guanaco}. Specifically, we measure the tendency of {Guanaco}-65B to generate socially biased outputs in comparison with alternative models (see Table~\ref{tbl:crows}). The mean bias score for {Guanaco}-65B is considerably lower relative to other base pretrained systems, suggesting that fine-tuning with the OASST1 dataset mitigates some biases inherent to the original LLaMA model. However, the extent to which {Guanaco} is robust against other types of biases remains unexamined, and comprehensive bias analysis for {Guanaco} and related conversational models is reserved for subsequent studies.

Another area not addressed in our experiments is the evaluation of varying bit-precisions, such as 3-bit representations, or alternative adapter techniques. In addition to LoRA, a multitude of Parameter Efficient FineTuning (PEFT) strategies have proven effective, though their scalability to very large models remains uncertain. While our selection of LoRA is supported by evidence of its reliability, it is possible that other adapters could further enhance performance. Given that post-quantization finetuning appears capable of restoring much of the information lost due to quantization, more aggressive quantization regimes may be feasible. For example, applying 3-bit GPTQ quantization to a base model followed by LoRA finetuning could potentially recover performance levels close to standard 16-bit full finetuning.

\section{Broader Impacts}

The proposed \textsc{QLoRA} method introduces, for the first time, the capacity to fine-tune models containing 33B parameters on consumer-grade GPUs and those with 65B parameters on a single professional GPU, achieving performance on par with traditional full-parameter finetuning approaches. Our strongest 33B-parameter model, trained on the Open Assistant dataset, demonstrates competitive results against ChatGPT in the Vicuna evaluation. Given that instruction finetuning is crucial for adapting pretrained large language models into effective chatbots resembling ChatGPT, we anticipate that our technique will democratize fine-tuning, especially benefiting researchers with restricted computational resources—thereby advancing the inclusivity and accessibility of cutting-edge NLP tools. As such, \textsc{QLoRA} offers a means to reduce the gap in technical capability between large corporate research labs and smaller academic or independent teams operating with limited hardware.

Another promising avenue for impact is the application of our method on mobile devices. We argue that \textsc{QLoRA} could represent a significant advancement by enabling finetuning of large language models (LLMs) directly on phones and other environments with limited computational resources. Although previous work has demonstrated that 7B parameter models can be executed on mobile phones, \textsc{QLoRA} is the first technique to facilitate finetuning of these models in such constrained settings. According to our estimates, it is possible to finetune approximately 3 million tokens overnight on an iPhone 12 Plus while charging. While the performance of finetuned 7B models does not yet match ChatGPT, we contend that the attained quality is sufficient to enable innovative uses previously constrained by concerns over privacy or inadequate LLM capabilities. By supporting on-device model ownership and control, \textsc{QLoRA} not only enhances user privacy but also simplifies the deployment process for LLMs.

Nevertheless, the capacity for finetuning carries inherent risks of misuse. While there are documented hazards associated with the proliferation of LLMs \citep{bommasani2021opportunities,bender2021dangers}, we maintain that democratizing access to this rapidly proliferating technology will promote more robust and independent evaluation than would be possible if such capabilities remain concentrated within large corporations that do not provide models or source code for public scrutiny.

In summary, we anticipate that \textsc{QLoRA} will substantially broaden access to high-quality LLM finetuning, thereby producing a net positive impact.